\bookchapter{خاتمة: ابقَ فضولياً، ابقَ يقظاً}

تهانينا على وصولك لهنا.

هذا الكتاب من عتاد الهاتف لمنظومة البرمجيات، من الإعدادات اليومية لوضع المطور، من الحماية الأمنية للتراث الرقمي. إن ساعدك في حل مشكلة عملية — حتى لو تعلمت إغلاق إذن غير ضروري، أو اكتسبت عادة نسخ احتياطي — فهذه الكلمات لم تُكتب عبثاً.

لكن أهم من أي مهارة محددة، أمران:

\begin{itemize}
  \item \textbf{ابقَ فضولياً}: الهاتف يتغير كل عام، AI يتغير، الشاشات القابلة للطي تتغير، الاتصال الفضائي يتغير. الخطوات التي تتعلمها اليوم قد تتقادم غداً، لكن فضول "أريد أن أعرف لماذا" لا يتقادم أبداً. أن تسأل "ما هذا الإذن" "أين تُخزن هذه البيانات" "هل تكشف هذه الوظيفة الخصوصية" أنفع من حفظ مئة خطوة.
  \item \textbf{ابقَ واعياً}: الهاتف أداة، ليس سيداً. صُمم ليجذب انتباهك ويشغل وقتك ويوجّه استهلاكك. استخدامه اختيارك، واستخدامك من قبله فقدان السيطرة. ارفع رأسك دورياً لشاشة الوقت، وفكّر "هل أستخدم الهاتف لإنجاز شيء، أم يستهلكني الهاتف؟" هذا ليس ترففاً، بل غريزة بقاء في العصر الرقمي.
\end{itemize}

التقنية تتغير، لكن احتياجات الإنسان لم تتغير — نحتاج الأمان والاتصال والاحترام. الهاتف مجرد ناقل. “درس الهاتف” الحقيقي ليس في الكتاب، بل في لحظة فتح قفل الشاشة كل يوم.

\begin{tcbpractice}{تمرين أخير}
أغلق هذا الكتاب، احمل هاتفك، افعل شيئاً واحداً:\textbf{فكّر 30 ثانية، ما الأشياء الثلاثة في هاتفك التي “فقدانها يؤلمك حقاً”؟}ثم تأكد من نسخها احتياطياً. هذا الكتاب يستحق.

\end{tcbpractice}
